\documentclass[12pt]{article}
\usepackage[margin=20mm]{geometry}
\usepackage{fontspec,amsmath,amssymb,booktabs,graphicx,hyperref}
\usepackage{polyglossia}
\usepackage{xcolor}
\ExplSyntaxOn
\NewDocumentCommand{\highlight}{m}{
  \group_begin:
  \setlength{\fboxsep}{0pt}
  \seq_set_split:Nnn \l_tmpa_seq {~} {#1}
  \seq_map_inline:Nn \l_tmpa_seq {\colorbox{yellow}{\strut ##1}~}
  \unskip
  \group_end:
}
\ExplSyntaxOff
\setdefaultlanguage{ukrainian}
\setotherlanguage{english}
\setmainfont{DejaVu Serif}
\newfontfamily\cyrillicfont{DejaVu Serif}
\setmonofont{DejaVu Sans Mono}
\newfontfamily\cyrillicfonttt{DejaVu Sans Mono}
\setlength{\emergencystretch}{4em}
\newcommand{\draft}[1]{\{#1\}}
\newcommand{\fig}[3]{\begin{figure}[htbp]\centering\includegraphics[width=.95\linewidth]{#1}\caption{#2}\label{#3}\end{figure}}
\title{Точність офсетів та обчислювальні компроміси в апроксимації полів знакової відстані компактними псевдогаусовими ядрами}

\begin{document}\maketitle
\begin{abstract}
    Поля знакової відстані використовують для неявного подання геометрії та побудови офсетів, однак збереження вибірки значень і відновлення поля можуть змінювати як самі відстані, так і геометрію множин рівня. Подання, узгоджене з нульовою множиною рівня, не обов'язково зберігає відстань поза нею. Суми компактних радіальних функцій забезпечують локальність внесків і свободу вибору точок інтерполяції. Це мотивує порівняння, у якому точність поля, градієнти, офсети, топологію та витрати ресурсів оцінюють окремо.

    Ми оцінили дев'ять реалізацій за зафіксованим обчислювальним протоколом \texttt{extended-4-6-n4-v1}. Параметри, дібрані на двох двовимірних геометріях, перенесено на шість незалежних геометрій. Медіанна середньоквадратична похибка (RMSE) поля в смузі для псевдогаусової інтерполяції у вузлах сітки становила 0.01353 у 2D та 0.01946 у 3D; для кубічної інтерполяції відповідні значення дорівнювали 0.01283 та 0.01847. Інтерполяція в точках на ізолініях не продемонструвала загальної переваги; хоч і її медіанна похибка у 3D становила 0.04961 \highlight{, проте з підібраними ізолініями під конкретну модель похибка становить 0.01????"}. Карта параметрів містила 13 конфігурацій у межах діагностичного порога, 19 із прапорцями якості та вісім із непокритими точками. Результати виявляють залежність компромісу від обчислювального навантаження в досліджених умовах; перенесення параметрів із 2D у 3D та дискретне оцінювання топології обмежують узагальнення.
\end{abstract}
\noindent\textbf{Ключові слова:} поле знакової відстані; компактна радіальна базисна функція; псевдогаусова інтерполяція; офсети; цифрова топологія; обчислювальні витрати.

\section{Постановка задачі}
Точне поле знакової відстані (SDF) задає відстань до межі, а знак указує, з якого боку межі розташована точка. У практичних обчисленнях інтерполянт замінює повторні звернення до точної геометрії. Його нульовий рівень може добре наближати межу, тоді як ненульові рівні відхиляються від точних офсетів. Тому задача подання стосується і геометричного положення, і точності відстаней, а також витрат на побудову, обчислення значень та зміну поля.

Ми досліджуємо компроміс, який забезпечує компактна псевдогаусова інтерполяція порівняно з ядрами Вендланда та сітковою інтерполяцією. Конкретне питання полягає в тому, чи поліпшує вибір точок інтерполяції на сусідніх рівнях SDF геометричну точність на незалежних прикладах за прийнятних обчислювальних витрат. Внесок роботи полягає у відтворюваному порівнянні й оцінюванні перенесення параметрів, а не у створенні нового ядра або доведенні загальної переваги.

\section{Огляд літератури}
...

\section{Методи}
\subsection{Побудова та позначення}
Ми використали аналітичні еталонні SDF та результати власних експериментів. Усі координати й відстані задано в безрозмірних одиницях порівняльного дослідження. Нехай $D\in\{2,3\}$, $\Omega=[-1,1]^D$, $h$ - крок сітки, $r$ - радіус підтримки та $\rho=r/h$. Центри задаються як $c_j=h\boldsymbol\ell$, де цілочисельні координати задовольняють $-K\le\ell_m\le K$, $K=\lceil(1+4h)/h\rceil$, а $j=1,\ldots,N$. Додатковий запас забезпечує наявність центрів поза областю оцінювання. Точки інтерполяції $y_i$, коефіцієнти $a_j$ і крок рівнів $\Delta$ мають різні ролі.

Для радіальної відстані $s\ge0$ зафіксоване ядро має вигляд
\begin{equation}\phi_{4,r}(s)=\left(1-s^2/r^2\right)_+^5,\qquad z_+=\max(z,0).\label{eq:kernel}\end{equation}
Нормування означає $\phi_{4,r}(0)=1$, а не одиничний інтеграл. Функція дорівнює нулю за $s\ge r$, а її похідні порядків від нульового до четвертого неперервно зануляються при $r$. Множник із нулем п'ятого порядку забезпечує гладкість $C^4$ на межі підтримки. Усередині підтримки радіальне продовження є поліномом від квадрата відстані. На рис.~\ref{fig:kernel} аналітично порівняно профілі; сама гладкість не доводить додатної визначеності чи точності.
\fig{figures/kernel.pdf}{Аналітичні профілі ядер залежно від $s/r$: псевдогаусове ядро з $n=4$ та нормовані формули Вендланда C2/C4. Інтерпольовану геометрію та емпіричні похибки не показано.}{fig:kernel}

Поле та систему для визначення його коефіцієнтів задають вирази
\begin{equation}
    f(x)=\sum_{j=1}^{N}a_j\phi_{4,r}(\|x-c_j\|),\qquad
    A_{ij}=\phi_{4,r}(\|y_i-c_j\|),\qquad Aa=d(y).
    \label{eq:fit}
\end{equation}
Тут $d(y)_i=d(y_i)$ для $i=1,\ldots,N$. Та сама конструкція застосовується у 2D та 3D зі зміною розмірності координат і евклідової норми. Система є квадратною, але з цього не випливає єдиність розв'язку. Після інтерполяції для обчислення поля потрібні коефіцієнти, центри або метадані для відновлення регулярної сітки, сімейство ядра, порядок і радіус. Точки інтерполяції не потрібні для обчислення значень, але їх зберігають для редагування з повторним використанням попередніх обчислень.

У сітковому варіанті $y_i=c_i$. У варіанті з ізолініями вибирають найближчу допустиму точку з об'єднання рівнів $d(y)=k\Delta$, $k\in\mathbb Z$, за умов $0<\Delta\le h$ і $\|y_i-c_i\|<r$. Реалізація перебирає рівні, для яких $|d(c_i)-k\Delta|\le r$, використовуючи 1-ліпшицевість точної SDF, та виконує геометричну проєкцію на кожен рівень. Порожні рівні повертають кандидати з нескінченними або невизначеними числовими значеннями, які відхиляються. Допускається зміщення менше за $r(1-10^{-12})$; вибір оновлюється лише за покращення понад $10^{-12}$. Рівні перебираються за зростанням, тому за однакової відстані зберігається менший рівень. Для неоднозначних геометричних проєкцій проєктори використовують детерміновані напрями. Неможливість знайти допустиму точку явно реєструється; точки, що збігаються після округлення до 11 десяткових знаків, відхиляються.

Система колокації розв'язується розрідженим LU-методом без регуляризації або збурення точок, що збігаються. Нескінченні чи невизначені коефіцієнти або максимальна вузлова нев'язка понад $10^{-7}\max(1,\max_i|d(y_i)|)$ спричиняють відмову. Розріджені запити відстаней cKDTree формують компактні взаємодії; у розширеному протоколі поле обчислюється пакетами по 1024 точки. На регулярній сітці кількість активних кандидатів становить $O((1+r/h)^D)$. Незалежність від загального числа центрів потребує фіксованого $D$, обмеженого $\rho$ та належного доступу до сусідів. Побудова дерев, пошук, факторизація та інші витрати реалізації додаються окремо; код будує дерева під час пакетного обчислення.

Поле використовується на $\Omega$. Поза об'єднанням скінченних підтримок воно дорівнює нулю, тому без правила продовження не є глобально точною SDF. Збіг на нульовому рівні не забезпечує $\|\nabla f\|=1$ або правильних ненульових офсетів. Наведені рівняння звірено з реалізацією; \draft{ПОТРІБНЕ ДЖЕРЕЛО: безпосередньо перевірити вбудовані рівняння наданого застарілого DOC-файлу щодо цієї постановки}. Допоміжні нотатки не вважаються дослівним відтворенням документа.

\subsection{Методи порівняння та зафіксований добір параметрів}
Для $u=s/r$ інші компактні профілі мають вигляд
\begin{equation}
    \psi_{C2}(u)=(1-u)_+^4(4u+1),\qquad
    \psi_{C4}(u)=(1-u)_+^6(35u^2+18u+3)/3.
\end{equation}
Кожен оцінюється з точками інтерполяції на сітці та на ізолініях. Лінійна інтерполяція й тензорно-добуткова кубічна схема not-a-knot у SciPy \textenglish{RegularGrid\-Interpolator} використовують ті самі центри та точність float64; кубічна схема застосовує розріджений прямий розв'язувач. Дев'ятий метод, cubic-plus-EDT, обчислює кубічний інтерполянт на сітці з кроком $h/2$, класифікує знак, віднімає внутрішнє евклідове перетворення відстані від зовнішнього та лінійно інтерполює отримані значення. Додатковий обсяг даних густішої сітки, попередня обробка й можливий дрейф поверхні є частиною порівняння.

Для добору використано коло з кутом $11^\circ$ і зсувом $-0.03$ та квадрат із кутом $23^\circ$ і зсувом $0.02$, $h=0.25$, початкові значення генератора (seed) 0--2 і по 512 рандомізованих точок Соболя на кожну вибірку. Цільовою функцією було середнє RMSE в смузі $|d|<0.25$ за шістьма комбінаціями форми та seed. Кандидат із хоча б однією невдалою вибіркою отримував нескінченне значення критерію. Для сіткових RBF використано $\rho\in\{1.25,1.5,2,2.5,3,4\}$; для ізоліній — $\rho\in\{1.5,2,3\}$ та $\Delta/h\in\{0.5,1\}$. Сіткові схеми без параметрів добору використовували одну конфігурацію. Перебрано всі кандидати й вибрано найменше скінченне значення критерію; це скінченний пошук, а не глобальний оптимум. Вибір записано до незалежного оцінювання.

У табл.~\ref{tab:selection} наведено зафіксовані параметри й загальний час добору, включно з невдалими кандидатами. Поле порядку зафіксовано на чотирьох для псевдогаусових методів, воно не впливає на сіткову інтерполяцію, а для Вендланда його роль визначає сімейство. Однакові обмеження кількості кандидатів не означають однакового обсягу геометричної інформації: звернення до SDF, виклики проєктора та кількість проєктованих точок обліковують різну роботу API; внутрішні операції проєктора не відображені в лічильниках SDF.
\begin{table}[htbp]\centering\small\begin{tabular}{lrrrr}\toprule
Метод & $\rho$ & $\Delta/h$ & Кандидати & Добір (с) \\ \midrule
pg\_grid & 2.0 & -- & 6 & 0.0847 \\
pg\_isolines & 2.0 & 0.5 & 6 & 0.1077 \\
wendland\_c2\_grid & 4.0 & -- & 6 & 0.0747 \\
wendland\_c2\_isolines & 3.0 & 0.5 & 6 & 0.1069 \\
wendland\_c4\_grid & 2.5 & -- & 6 & 0.0752 \\
wendland\_c4\_isolines & 3.0 & 0.5 & 6 & 0.1075 \\
linear & -- & -- & 1 & 0.0020 \\
cubic & -- & -- & 1 & 0.0067 \\
cubic\_redistance & -- & -- & 1 & 0.0079 \\
\bottomrule\end{tabular}
\caption{Зафіксований добір і збережені витрати часу. Крок рівнів не застосовується до інтерполяції у вузлах сітки, а радіус — до сіткових інтерполяційних схем.}\label{tab:selection}\end{table}

Контрольні геометрії охоплювали капсулу, фігуру із заокругленими кутами, багатокутник із вузьким перешийком, два диски, сферу й тор. Seed 10--12 задавали нові кути, рівномірно розподілені в $[31,68]$ градусів, і зсуви з $[-0.08,0.08]$. Використано $h=0.2,0.3$ у 2D та $0.4,0.5$ у 3D, офсети $t=-0.1,0,0.1$ і сітки виділення рівнів $81^2$ та $31^3$. Параметри, дібрані у 2D, безпосередньо перенесено у 3D. Отримано 234 навчальні оцінки та 972 записи офсетних метрик на контрольних геометріях; офсети одного поля не є незалежними експериментами.

\subsection{Точність, топологія та ресурси}
Для множини точок сітки виділення рівнів $P$ позначимо $B_t=\{x\in P:|d(x)-t|<0.2\}$. Похибка поля у смузі та дискретна відстань між офсетами визначаються як
\begin{equation}
    E_t=\sqrt{|B_t|^{-1}\sum_{x\in B_t}(f(x)-d(x))^2},\quad
    H_t=\max\{\max_{v\in V_f}\min_{w\in V_d}\|v-w\|,\max_{w\in V_d}\min_{v\in V_f}\|w-v\|\}.
\end{equation}
$V_f,V_d$ — вершини, отримані виділенням рівня $t$. $H_t$ є відстанню між хмарами точок, а не гарантованою відстанню Гаусдорфа між неперервними поверхнями. Порожні або нерозділені рівні реєструються, їм не приписується нульова похибка. Значення точної SDF у вершинах модельної поверхні є окремою геометричною діагностикою.

Вимірювання градієнтів використовують 2048 точок Соболя з $|d|<0.25$. Еталонні центральні різниці з кроками $2\times10^{-5}$ і $10^{-5}$ приймаються, якщо норма різниці їхніх векторів менша за $10^{-4}$, а норма градієнта, обчисленого з меншим кроком, відрізняється від одиниці менш ніж на $10^{-4}$. Ця чисельна маска не характеризує всю медіальну вісь. Градієнти компактних полів обчислюються аналітично; для сіткових полів застосовано центральні різниці з кроком $h\times10^{-5}$, де для EDT береться крок густішої сітки. RMSE градієнта дорівнює $\sqrt{\operatorname{mean}\|g-g_d\|^2}$, а RMSE ейкональної нев'язки — $\sqrt{\operatorname{mean}(\|g\|-1)^2}$. Узгодженість нормалей — це середнє $g\cdot g_d/(\|g\|+10^{-30})$ на прийнятій множині, де $g_d$ чисельно має одиничну довжину.

Цифрові тіла задаються умовами $f\le t$ і $d\le t$. Зв'язність об'єкта дорівнює 4/6, фону — 8/26. Обчислюються кількість компонент $b_0$, кількість порожнин $b_2$ у 3D та ейлерова характеристика $\chi$; звідси $b_1=b_0-\chi$ у 2D або $b_1=b_0+b_2-\chi$ у 3D. Контакт із межею області реєструється окремо. У дослідженні роздільності поля залишаються незмінними: seed 42, $h=0.2$ у 2D і $0.4$ у 3D, сітки $51^2,101^2,201^2$ та $25^3,41^3,65^3$. Офсети становлять $(-0.1,0,0.1)$ для заокругленої фігури та сфери, $(0.06,0.08,0.10)$ для дисків, $(-0.12,-0.10,-0.08)$ для перешийка та $(-0.19,0,0.30,0.34)$ для тора. Ці 432 записи є емпіричною перевіркою впливу роздільності, а не доведенням збіжності.

Карта із 40 випадків варіює $h\in\{0.2,0.3\}$, $\rho\in\{0.5,1,1.5,2,3\}$, $\Delta/h\in\{0.5,1\}$ та стратегію за $n=4$, використовуючи окремий запас $3h$. Умова $H>0.1$ або невідповідність цифрової топології задає діагностичний прапорець на сітці $101^2$. Сіткова стратегія не використовує крок рівнів, тому деякі конфігурації карти повторюють те саме поле.

Ресурсні вимірювання охоплюють 36 випадків в ізольованих підпроцесах, один потік BLAS/OMP, три повтори та пакети з 256/2048/8192 запитів після прогріву. Геометрії — капсула у 2D та тор у 3D із кутом $47^\circ$ і зсувом $0.01$. Обсяг числових даних поля (payload) охоплює числові масиви; метадані, збережені LU-фактори, пікову резидентну пам'ять процесу (RSS), час побудови та проєктування розрізняють. Абсолютний піковий RSS включає імпорт бібліотек, а піковий RSS редагування є накопиченим максимумом, до якого входять попередні оцінювання якості. Редагування з кешуванням повторно використовує початкові точки інтерполяції та факторизацію; за повної перебудови точки вибираються знову. Отже, кешований режим не здійснює адаптивного перевибору точок на ізолініях. Домінування за Парето визначається окремо для кожної розмірності: не гірші середня похибка офсетів, обсяг даних і час запитів та строге покращення хоча б одного з них. Однакові вектори критеріїв не домінують один над одним.

\section{Вибір точок інтерполяції та порівняльне оцінювання}
У табл.~\ref{tab:errors} наведено медіани за успішними контрольними записами та діапазони трьох медіан, обчислених окремо для кожного seed. Скрипт підготовки звіряє їх із первинними JSON і збереженим підсумком. Один із 36 тривимірних записів псевдогаусового методу з ізолініями не мав визначеної відстані між офсетами; його виключено лише з медіани $H$. Усі інші групи метод/розмірність мали визначені відстані в кожному записі. Медіани сіткового псевдогаусового методу були близькими до кубічної інтерполяції, що не доводить еквівалентності. Найменшу медіанну RMSE поля в смузі у 2D мав сітковий метод Вендланда C2, у 3D — сітковий метод Вендланда C4. Ранжування стосується саме цих форм, кроків і бюджетів кандидатів.
\begin{table}[htbp]\centering\resizebox{\linewidth}{!}{\begin{tabular}{lrrrrr}\toprule
Метод & $D$ & RMSE у смузі & Діапазон медіан за seed & $H$ & RMSE градієнта \\ \midrule
cubic & 2 & 0.012831 & 0.012555--0.013453 & 0.038722 & 0.22797 \\
cubic\_redistance & 2 & 0.060900 & 0.059132--0.065543 & 0.104718 & 0.72484 \\
linear & 2 & 0.022883 & 0.022656--0.023188 & 0.085709 & 0.33686 \\
pg\_grid & 2 & 0.013528 & 0.013100--0.014739 & 0.037584 & 0.23008 \\
pg\_isolines & 2 & 0.014654 & 0.014542--0.015261 & 0.042950 & 0.24225 \\
wendland\_c2\_grid & 2 & 0.012777 & 0.012757--0.013526 & 0.039340 & 0.22751 \\
wendland\_c2\_isolines & 2 & 0.013766 & 0.013043--0.014271 & 0.041500 & 0.23963 \\
wendland\_c4\_grid & 2 & 0.013254 & 0.012420--0.013313 & 0.041773 & 0.22887 \\
wendland\_c4\_isolines & 2 & 0.013778 & 0.013414--0.014597 & 0.042325 & 0.23988 \\
cubic & 3 & 0.018466 & 0.017940--0.018847 & 0.046052 & 0.18482 \\
cubic\_redistance & 3 & 0.105771 & 0.102410--0.108088 & 0.148238 & 0.93009 \\
linear & 3 & 0.067278 & 0.066975--0.068381 & 0.136773 & 0.34608 \\
pg\_grid & 3 & 0.019462 & 0.018871--0.019830 & 0.046579 & 0.21465 \\
pg\_isolines & 3 & 0.049607 & 0.042961--0.054275 & 0.163485 & 0.47347 \\
wendland\_c2\_grid & 3 & 0.017688 & 0.017577--0.018002 & 0.047756 & 0.17112 \\
wendland\_c2\_isolines & 3 & 0.022730 & 0.021852--0.023964 & 0.053608 & 0.18305 \\
wendland\_c4\_grid & 3 & 0.015525 & 0.015525--0.016084 & 0.041024 & 0.17011 \\
wendland\_c4\_isolines & 3 & 0.026734 & 0.025780--0.027266 & 0.060094 & 0.20170 \\
\bottomrule\end{tabular}
}\caption{Описові медіани на контрольних геометріях. Діапазони за seed не є довірчими інтервалами; $H$ — дискретна відстань між офсетами лише для визначених рівнів. Кількості невизначених відстаней збережено в записі походження даних.}\label{tab:errors}\end{table}

Зміна точок інтерполяції зберігає ті самі центри та кількість коефіцієнтів, але змінює обмеження й підготовчі обчислення. На рис.~\ref{fig:good} та \ref{fig:bad} відтворено фактичні контрольні пари відповідно з найменшим і найбільшим двовимірним відношенням похибки поля в смузі для ізоліній до сіткового варіанта. Це цілеспрямовані ілюстрації сприятливого й несприятливого випадків, а не репрезентативна випадкова вибірка. Їхні виміряні похибки повторно звірено зі збереженими записами. Різноспрямовані результати не встановлюють причинного пояснення відмінностей у точності.
\fig{figures/improvement.pdf}{Капсула; seed 11, $h=0.3$, $\rho=2.0$, $\Delta/h=0.5$, офсет 0.0, роздільність $81^2$. RMSE у смузі для сітки/ізоліній: 0.018269/0.014569. Пунктир — еталон; синій колір — зображене поле.}{fig:good}
\fig{figures/failure.pdf}{Два кола; seed 12, $h=0.3$, $\rho=2.0$, $\Delta/h=0.5$, офсет $-0.1$, роздільність $81^2$. RMSE у смузі для сітки/ізоліній: 0.016207/0.026605. Пунктир — еталон; синій колір — зображене поле.}{fig:bad}

Зафіксований добір вибрав $\rho=2$ для обох псевдогаусових варіантів і $\Delta/h=0.5$ для ізоліній. Медіанна RMSE варіанта з ізолініями у 3D, 0.049607, значно перевищувала сіткове значення 0.019462. Оскільки добір використовував лише двовимірні форми, це результат спостереження за перенесенням параметрів; він не показує, чого можна досягти за окремого добору у 3D. Уточнення на навчальних формах і оцінювання на контрольних відповідають на різні питання.

\section{Обговорення}
Карта параметрів містила 13 конфігурацій у межах порога, 19 із прапорцями якості та вісім із непокритими точками. Ці частки не є ймовірностями успіху методу. Мала вузлова нев'язка може співіснувати з неточними офсетами або незв'язними дискретними тілами, оскільки інтерполяція обмежує лише вибрані значення. Радіус підтримки змінює дальність взаємодії, крок сітки — дискретизацію, а крок рівнів — умови інтерполяції. Гладкість у цьому дослідженні зафіксована; дані не виокремлюють вплив зміни її порядку.

Випадок тонкого тора конкретно демонструє обмеження дискретизації (табл.~\ref{tab:topology}). За офсету $-0.19$ початковий радіус трубки $0.20$ стає рівним $0.01$, що мало порівняно з кроком сітки виділення поверхні. Навіть точний еталон може мати неправильну кількість цифрових компонент або отворів. Порівняння цифрової топології моделі й еталона корисне, але їхній збіг не доводить правильності неперервної топології; прапорець контакту з межею також обмежує інтерпретацію повного тіла.
\begin{table}[htbp]\centering\small\begin{tabular}{lrrr}\toprule
Метод & Роздільність & Еталон $(b_0,b_1,b_2)$ & Модель $(b_0,b_1,b_2)$ \\ \midrule
pg\_grid & 25 & (10, 0, 0) & (0, 0, 0) \\
pg\_grid & 41 & (18, 0, 0) & (0, 0, 0) \\
pg\_grid & 65 & (39, 0, 0) & (0, 0, 0) \\
pg\_isolines & 25 & (10, 0, 0) & (4, 0, 0) \\
pg\_isolines & 41 & (18, 0, 0) & (4, 0, 0) \\
pg\_isolines & 65 & (39, 0, 0) & (4, 0, 0) \\
cubic & 25 & (10, 0, 0) & (0, 0, 0) \\
cubic & 41 & (18, 0, 0) & (0, 0, 0) \\
cubic & 65 & (39, 0, 0) & (0, 0, 0) \\
\bottomrule\end{tabular}
\caption{Тонкий тор, офсет $-0.19$, seed 42, фіксований $h=0.4$; PG: $n=4$, $\rho=2$, для ізоліній $\Delta/h=0.5$. Роздільність змінюється за незмінних полів.}\label{tab:topology}\end{table}

Ресурсний приклад тора в табл.~\ref{tab:resources} показує, чому самого ранжування за похибкою недостатньо. Для всіх методів використано ту саму збережену геометрію та $h=0.4$; дискретні похибки офсетів усереднюються за рівнями $-0.1,0,0.1$. Час запитів є медіаною для 8192 запитів; добір параметрів оплачується додатково. Результати характеризують ці реалізації та це навантаження, а не всі контрольні форми чи промислову продуктивність. Відомостей про обладнання, окрім збережених даних про програмне забезпечення та платформу, недостатньо для порівняння з урахуванням апаратних відмінностей.
\begin{table}[htbp]\centering\resizebox{\linewidth}{!}{\begin{tabular}{lrrrrr}\toprule
Метод & $h$ & Дані (байт) & Побудова (с) & Запити (с) & $H$ \\ \midrule
pg\_grid & 0.4 & 27000 & 0.0637 & 0.0220 & 0.05987 \\
pg\_isolines & 0.4 & 27000 & 0.0841 & 0.0218 & 0.11449 \\
wendland\_c2\_grid & 0.4 & 27000 & 0.8532 & 0.0979 & 0.06914 \\
wendland\_c2\_isolines & 0.4 & 27000 & 0.3845 & 0.0478 & 0.09767 \\
wendland\_c4\_grid & 0.4 & 27000 & 0.1859 & 0.0332 & 0.07421 \\
wendland\_c4\_isolines & 0.4 & 27000 & 0.3892 & 0.0474 & 0.09455 \\
linear & 0.4 & 27000 & 0.0004 & 0.0006 & 0.18607 \\
cubic & 0.4 & 27000 & 0.0904 & 0.0025 & 0.06634 \\
cubic\_redistance & 0.4 & 195112 & 0.1009 & 0.0006 & 0.14481 \\
\bottomrule\end{tabular}
}\caption{Збережені ресурсні випадки для тора: $h=0.4$, кут $47^\circ$, зсув $0.01$, сітка виділення поверхні $31^3$; параметри з табл.~\ref{tab:selection}.}\label{tab:resources}\end{table}

Кешовані оновлення та повні перебудови можуть відрізнятися і точністю, і часом (табл.~\ref{tab:editing}). Наведені похибки — це медіани для трьох змінених положень на 256 точках Соболя, а не контрольна метрика поля в смузі. Збереження факторизації потребує пам'яті понад числові дані поля; для оновлення з фіксованими вузлами також потрібні збережені точки. Тому навантаження з великою кількістю запитів або редагувань може надавати перевагу іншій конфігурації. Окремі критерії Парето не враховують побудови та добору, тому ці витрати слід враховувати під час вибору практичної реалізації.
\begin{table}[htbp]\centering\resizebox{\linewidth}{!}{\begin{tabular}{lrrrrr}\toprule
Метод & Підгот. кешу (с) & Оновлення (с) & Перебудова (с) & RMSE кеш. & RMSE перебуд. \\ \midrule
pg\_grid & 0.0655 & 0.00121 & 0.0638 & 0.01695 & 0.01695 \\
pg\_isolines & 0.0853 & 0.00123 & 0.0846 & 0.02378 & 0.02442 \\
wendland\_c2\_grid & 0.8615 & 0.00392 & 0.8472 & 0.01267 & 0.01267 \\
wendland\_c2\_isolines & 0.3860 & 0.00251 & 0.5097 & 0.01953 & 0.02108 \\
wendland\_c4\_grid & 0.1991 & 0.00182 & 0.1841 & 0.01178 & 0.01178 \\
wendland\_c4\_isolines & 0.3946 & 0.00240 & 0.5163 & 0.01854 & 0.02009 \\
linear & 0.0000 & 0.00014 & 0.0002 & 0.03844 & 0.03844 \\
cubic & 0.1553 & 0.00219 & 0.0892 & 0.01205 & 0.01205 \\
cubic\_redistance & 0.1533 & 0.01061 & 0.0982 & 0.07265 & 0.07265 \\
\bottomrule\end{tabular}
}\caption{Редагування для ресурсних випадків тора з табл.~\ref{tab:resources}. RMSE — медіана збережених похибок для трьох редагувань.}\label{tab:editing}\end{table}

Порівняльне дослідження забезпечує кількісне оцінювання на скінченному наборі аналітичних геометрій. Воно не встановлює загальної збіжності, глобальної стійкості, промислової продуктивності або унікальності щодо літератури. Окремий добір у 3D, зовнішні геометрії, більші масштаби та обґрунтовані оцінки похибок залишаються майбутньою роботою. Архівні тести перевіряють реалізацію та властивості протоколу, а не математичну перевагу.

\section{Висновки}
У роботі порівняно дев'ять реалізацій за зафіксованим протоколом. Псевдогаусова інтерполяція у вузлах сітки забезпечила точність поля в смузі, близьку до кубічної інтерполяції, у досліджених умовах. Вибір точок на ізолініях дав покращення в окремих конфігураціях, що вказує на потенціал цього підходу за відповідного добору параметрів; водночас отримані результати не встановлюють його загальної переваги. Внесок полягає у відтворюваному порівняльному оцінюванні з явними обмеженнями перенесення параметрів.

Компактність підтримки обмежує кількість активних кандидатів за фіксованих розмірності та $r/h$. Повні обчислювальні витрати також включають операції з деревами, побудову та розріджену факторизацію. Збережені часи запитів і підготовки слід інтерпретувати окремо.

Компактна сума зберігає по одному коефіцієнту для кожного регулярного центра разом із метаданими сітки та ядра. Обсяг коефіцієнтів може дорівнювати розміру масиву сіткових значень за однакової точності, але це не означає рівності загального обсягу зберігання, кешованих факторів чи пікової пам'яті побудови.

Свобода вибору точок інтерполяції справді існує та дала локальні покращення в збережених порівняннях. Вона також спричинила погіршення, зокрема зростання сукупної похибки у 3D. Отже, належне правило вибору точок залежить від геометричних вимог і обчислювального навантаження.

Надана реалізація на Python/SciPy обчислює поля через пакетні компактні взаємодії. Її метадані та коефіцієнти забезпечують відтворюване відновлення; сумісність із зовнішніми промисловими форматами цим дослідженням не встановлена.

Набори кандидатів залишають невирішеними питання взаємодії параметрів, особливо у 3D. Незалежний добір має оцінювати поле, градієнт, офсети й топологію разом із покриттям і ресурсами, а не оптимізувати одну агреговану похибку без перевірки наслідків.

У подальших дослідженнях потрібно окремо підібрати параметри для 3D, перевірити метод на нових геометріях та обґрунтувати його стійкість і збіжність. Це покаже, чи справджуються отримані висновки для інших форм, окрім досліджених аналітичних прикладів.

\begin{english}
    \renewcommand{\refname}{\textukrainian{Список літератури}}
    \bibliographystyle{unsrt}\bibliography{references}
\end{english}
\end{document}
